% BOOK-FILE-ID: f-5389cc593596
\subsection{対称式を使う}\label{節：対称式を使う}
導入で紹介した3つ目の\bookterm{recurrence}{漸化式}を考える.
\[
  a_1=3, \quad a_{n+1}=a_n^2-2
\]
邪魔なのは$-2$である.これをうまく消すには,二乗を展開したときに$2$が生まれてくれればよい.
そこで思い出したいのが,
\[
  (a+b)^2=a^2+2ab+b^2
\]
である.この$2ab$に注目すると,\,$ab=1$ならちょうど$2$が生まれる.
そこで積が$1$になる2数として,\,$a=\alpha,\ b=\frac{1}{\alpha}\ (\alpha\ne0)$と選んでやる.
\[
  \left(\alpha+\frac{1}{\alpha}\right)^2=\alpha^2+2+\frac{1}{\alpha^2}
\]
だから,\,$\alpha+\frac1\alpha$を2乗して2を引くと,\,$\alpha^2+\frac{1}{\alpha^2}$になるのだ.
つまり,\,$a_n=\alpha^{k}+\alpha^{-k}$の形をしていれば,1回進むごとに指数$k$が2倍になる.
第$n$項の指数は$2^{n-1}$であり,あとは初項に合う$\alpha$を選べばよい.

\begin{bookpoint}
  \textbf{POINT}\par
  $a_n=\alpha^{k}+\alpha^{-k}$と置いて,漸化式が指数の何倍にあたるかを読む.
\end{bookpoint}


\noindent \bookblocklink{example-symmetric-1}{problem}{answer}{【問題】}\par\nobreak
\begin{smallquote}
  次の\bookterm{recurrence}{漸化式}で定義される\bookterm{sequence}{数列}の\bookterm{general}{一般項}を求めよ.
  \[
    a_1=3, \quad a_{n+1}=a_n^2-2
  \]
\end{smallquote}
\bookblocklink{example-symmetric-1}{answer}{problem}{【解答】}
\begin{smallquote}
  \bookterm{general}{一般項}が
  \[
    b_n=\alpha^{2^{n-1}}+\alpha^{-2^{n-1}}
  \]
  で表される\bookterm{sequence}{数列}$b_n$を考える.
  \begin{align*}
    b_{n+1}&=\alpha^{2^{n}}+\alpha^{-2^{n}}\\
    &=\left(\alpha^{2^{n-1}}+\alpha^{-2^{n-1}}\right)^2-2\\
    &=b_n^2-2
  \end{align*}
  と漸化式が一致する.初項を合わせると,
  \begin{align*}
    &\alpha+\alpha^{-1}=3\\
    &\alpha^2-3\alpha+1=0\\
    &\therefore \quad \alpha=\frac{3\pm \sqrt{5}}{2}
  \end{align*}
  2つの解は互いに逆数で,同じ$b_n$を与える.初項と漸化式の一致から,
  \[
    a_n=\left(\frac{3+ \sqrt{5}}{2}\right)^{2^{n-1}}+\left(\frac{3- \sqrt{5}}{2}\right)^{2^{n-1}}
  \]
\end{smallquote}

ここで,定数を $-2$ から $-1$ に変えてみよう.
\[
  a_{n+1}=a_n^2-1
\]
同じ置換では,
\begin{align*}
  a_n^2-1
  &=\left(u+u^{-1}\right)^2-1\\
  &=u^2+u^{-2}+1,
\end{align*}
と余分な1が残り,同じ形で二乗を繰り返すことができない.
係数が似ているだけではだめで,恒等式とぴったり一致していることが必要なのである.

同じ考え方は,3次式にも使える.
\[
  \left(\alpha+\frac{1}{\alpha}\right)^3
  -3\left(\alpha+\frac{1}{\alpha}\right)
  =\alpha^3+\frac{1}{\alpha^3}
\]
対称な組み合わせを保ったまま,指数を3倍できる.

和を取るか差を取るかで,打ち消せる項の符号も変わる.
差の三乗を展開すれば,
\[
  \left(\alpha-\frac{1}{\alpha}\right)^3
  +3\left(\alpha-\frac{1}{\alpha}\right)
  =\alpha^3-\frac{1}{\alpha^3}
\]
となる.指数$3^{n-1}$の候補を例題で確かめよう.

\noindent\bookblocklink{example-symmetric-2}{problem}{answer}{【問題】}\par\nobreak
\begin{smallquote}
  $a_1=\frac{8}{3},\quad a_{n+1}=a_n^3+3a_n$
\end{smallquote}
\bookblocklink{example-symmetric-2}{answer}{problem}{【解答】}
\begin{smallquote}
  \bookterm{general}{一般項}が
  \[
    b_n=\alpha^{3^{n-1}}-\alpha^{-3^{n-1}}
  \]
  で表される\bookterm{sequence}{数列}$b_n$を考える.
  \begin{align*}
    b_{n+1}&=\alpha^{3^{n}}-\alpha^{-3^{n}}\\
    &=b_n^3+3b_n
  \end{align*}
  と漸化式が一致する.初項を合わせると,
  \begin{align*}
    &\alpha-\alpha^{-1}=\frac{8}{3}\\
    &3\alpha^2-8\alpha-3=0\\
    &\therefore \quad \alpha=3,\,-\frac13
  \end{align*}
  $3^{n-1}$は奇数なので,どちらを選んでも$b_n$は同じ値になる.
  初項と漸化式の一致から,\,$\alpha=3$を選べば,
  \[
    a_n=3^{3^{n-1}}-3^{-3^{n-1}}
  \]
\end{smallquote}
